\subsection{The limiting Poisson optimizer}
\label{sec:atom-count-law}

To prove Theorem~\ref{thm:atom-count-limit}, we must identify a canonical
limit for the compact likelihoods used above. Replacing the moving variance
penalty by its hard wall produces the optimizer \(\nu_b^*\) from
\eqref{eq:atom-limit-functional}; its uniqueness turns support tightness into
convergence to a law depending only on \(r\).

The functional \(\cJ_{N,b}\) agrees with \(\mathcal J_N^K\) in
\eqref{eq:support-one-d-limit-functional} whenever \(\nu\) is supported on
the compact interval \(K\subset[b,1)\). Its first variation is
\begin{equation}
\label{eq:atom-limit-kkt}
    \Gamma_\nu(\theta)
    =S(\theta)-\int_{\mathbb R}
       e^{\theta x}\frac{F_\nu(x)}{1+F_\nu(x)}\,N(dx).
\end{equation}
Thus a maximizing measure satisfies \(\Gamma_\nu\le0\) on \([b,1)\),
with equality on its support.

The negative Poisson points form a uniqueness set for the Laplace transforms
appearing in \eqref{eq:atom-limit-functional}. This is much stronger than
uniqueness on any fixed finite set of observations.
\begin{lemma}
\label{lem:atom-uniqueness-set}
Write the negative points of \(N\) as \(-T_j\). Almost surely, if a finite
signed measure \(\delta\) is supported on a compact interval
\([b,c]\subset(0,\infty)\) and
\[
    \int_b^c e^{-T_j\theta}\,\delta(d\theta)=0
    \qquad\text{for every }j,
\]
then \(\delta=0\).
\end{lemma}

\begin{proof}
For \(\operatorname{Re}z>0\), set
\[
    H(z)=e^{bz}\int_b^c e^{-z\theta}\,\delta(d\theta).
\]
This is a bounded analytic function on the right half-plane. Unless it
vanishes identically, its zeros satisfy the half-plane Blaschke condition
\cite[Chapter~II]{garnett2007bounded},
\[
    \sum_{H(z)=0}\frac{\operatorname{Re}z}{1+|z|^2}<\infty,
\]
with multiplicity.

The points \(T_j\) form a Poisson process on \((0,\infty)\) with intensity
\(e^t\,dt\). If \(Q_m\) counts its points in \([m,m+1]\), then the
\(Q_m\)'s are independent Poisson variables with means \((e-1)e^m\).
A Chernoff bound and Borel--Cantelli give
\(Q_m\ge(e-1)e^m/2\) for all sufficiently large \(m\), almost surely.
Consequently,
\[
    \sum_j\frac{T_j}{1+T_j^2}
    \ge\sum_{m\ \mathrm{large}}
       \frac{mQ_m}{1+(m+1)^2}=\infty.
\]
Thus \(H(T_j)=0\) for every \(j\) forces \(H\equiv0\). Uniqueness of the
Laplace transform then gives \(\delta=0\).
\end{proof}


We can now attach a canonical optimizer to each hard-wall location. Its
support size will be the one-sided contribution to the limiting atom count.
\begin{proposition}
\label{prop:atom-limit-unique}
For every fixed \(b\in(1/2,1)\), the functional \(\cJ_{N,b}\) has an almost
surely unique maximizer \(\nu_b^*\). This measure has finite support in a
random compact subinterval of \([b,1)\).
\end{proposition}

\begin{proof}
Since \(S(\theta)\to-\infty\) as \(\theta\uparrow1\), there is a random
\(c_N\in(b,1)\) such that \(S(\theta)<0\) on \([c_N,1)\). The second term in
\eqref{eq:atom-limit-kkt} is nonnegative, so deleting any mass in this
interval strictly increases the objective. Proposition
\ref{prop:support-one-d-analytic-active-set}, applied on a compact interval
containing \([b,c_N]\), gives existence and finite support.

The first term in \eqref{eq:atom-limit-functional} is linear, while \(h\) is
strictly concave. If two measures \(\nu_0,\nu_1\) were maximizers, equality
in the concavity inequality for their midpoint would require
\(F_{\nu_0}(X)=F_{\nu_1}(X)\) at every point \(X\) of \(N\). In particular,
the two transforms agree at every negative point. Lemma
\ref{lem:atom-uniqueness-set}, applied to \(\nu_1-\nu_0\), gives
\(\nu_0=\nu_1\).
\end{proof}

For later use, write
\[
    K_b=|\supp(\nu_b^*)|.
\]
The value zero is allowed: \(K_b=0\) exactly when the one-sided Poisson field
does not cross zero on \([b,1)\).

\subsection{The moving penalty and the hard wall}

The next step in Theorem~\ref{thm:atom-count-limit} shows why no
subpolynomial feature of \(\eps_n=n^{-r_n}\) survives in the limiting
optimizer. The deterministic penalty changes abruptly at
\(b_r=1-\sqrt r\), but its value exactly at that point need not converge.
A variational argument moves any boundary mass an asymptotically negligible
distance to the vanishing-penalty side.

Recall \(\lambda_n=\eps_nu_n^2/\{2(1-\eps_n)\}\), the scale
\(c_n(\theta)\) from Section~\ref{sec:one-dimensional-edge-field}, and the
penalty
\[
    \cP_n(\theta)
    =\frac{1-e^{-\lambda_n\theta^2}}{c_n(\theta)}
\]
from \eqref{eq:support-deterministic-penalty}. Assume
\(r_n\to r\in(0,1/4)\). Since
\(u_n^2/2=\log n+O(\log\log n)\) and \(\lambda_n\to0\), a direct expansion
gives, uniformly on compact subintervals of \((1/2,1)\),
\begin{equation}
\label{eq:atom-penalty-exponent}
    \log\cP_n(\theta)
    =\{(1-\theta)^2-r_n\}\log n+O(\log\log n).
\end{equation}
Consequently, for every fixed \(\eta>0\),
\begin{equation}
\label{eq:atom-hard-wall}
    \sup_{\theta\ge b_r+\eta}\cP_n(\theta)\longrightarrow0,
    \qquad
    \inf_{\theta\le b_r-\eta}\cP_n(\theta)\longrightarrow\infty.
\end{equation}

The following deterministic lemma records the variational consequence. It is
stated on a compact interval because all random likelihoods have already been
localized there.
\begin{lemma}
\label{lem:atom-hard-wall}
Let \(K=[a,c]\Subset(1/2,1)\), with \(a<b_r<c\), and suppose concave
functionals \(\widetilde L_n\) converge locally uniformly on every
bounded-mass set of finite positive measures on \(K\) to
\(\mathcal J_N^K\). Define
\[
    L_n(\nu)=\widetilde L_n(\nu)-\int_K\cP_n(\theta)\,\nu(d\theta).
\]
If the masses of the maximizing measures are tight, every weak limit of
maximizers of \(L_n\) is supported on \([b_r,c]\) and maximizes
\(\mathcal J_N^K\) there. Conversely, every measure on \([b_r,c]\) has a
recovery sequence with the same limiting value.
\end{lemma}

\begin{proof}
The two bounds in \eqref{eq:atom-hard-wall} make the penalty vanish uniformly
on compact subsets of \((b_r,c]\) and force every bounded-objective sequence
to put asymptotically no mass on \([a,b_r-\eta]\), for each \(\eta>0\).
Tightness of the masses gives the asserted support and the usual limsup
inequality.

For recovery, put
\[
    \delta_n=|r_n-r|^{1/2}
       +\left(\frac{\log\log n}{\log n}\right)^{1/2}.
\]
Then \(\delta_n\downarrow0\) and
\(\cP_n(b_r+\delta_n)\to0\) by
\eqref{eq:atom-penalty-exponent}. Leave mass above
\(b_r+\delta_n\) fixed and move all mass in
\([b_r,b_r+\delta_n]\) to \(b_r+\delta_n\). The penalty is decreasing on
a fixed neighborhood of \(b_r\) for all large \(n\), as follows by
differentiating its definition, so the recovery penalties vanish. Local
uniform convergence of \(\widetilde L_n\) completes the proof.
\end{proof}


The compact expansion already used for support tightness becomes uniform
across the moving wall after the deterministic penalty is added back. We
record the strengthened form needed to identify the optimizer itself.
\begin{lemma}
\label{lem:atom-wall-functional-convergence}
Let \(K=[a,c]\Subset(1/2,1)\) contain \(b_r\), and assume
\[
    \frac r2+(1-a)^2<\frac12.
\]
On every set of edge-measure pairs with total mass at most \(M\), the
functionals
\[
    \mathcal L_{n,K}(\nu_+,\nu_-)
    +\int_K\cP_n\,d\nu_++\int_K\cP_n\,d\nu_-
\]
converge uniformly to
\(\mathcal J_{N_+}^K(\nu_+)+\mathcal J_{N_-}^K(\nu_-)\). The first
variations converge jointly and locally uniformly, together with their first
two derivatives, on every smaller complex neighborhood. Here \(N_+\) and
\(N_-\) are independent copies of \(N\).
\end{lemma}

\begin{proof}
Truncate the two Gaussian edge processes to a bounded edge window. On this
window the kernel and its first two parameter derivatives are uniformly
bounded and equicontinuous over measures of mass at most \(M\), so compact
point-process convergence gives the assertion. On the negative tail, the
nonlinear likelihood remainder has envelope \(C_Me^{2ax}\); its
differentiated first variations are bounded by a polynomial in \(|x|\) times
\(C_Me^{(a+\operatorname{Re}\theta)x}\). These envelopes are integrable
because \(2a>1\). The positive tail contains only finitely many limiting
points. A finite weak net of the compact measure set makes the truncation
uniform.

The deterministic part of the normalized first variation is exactly
\(-\cP_n\). Once it is added back,
\(e^{-\lambda_n\theta^2}=1+o(1)\) uniformly with two derivatives on \(K\),
so the nonlinear limit is unchanged. Translation by the central atom is
negligible with two derivatives under the displayed condition on \(a\).
Finally, integrate the uniform first-variation convergence along
\(s\mapsto(s\nu_+,s\nu_-)\), \(0\le s\le1\), to obtain convergence of the
functionals. Joint convergence of the two extreme point processes gives
independent \(N_+\) and \(N_-\).
\end{proof}

An atom of physical mass \(w\) at signed score displacement
\(\theta u_n\) from the central atom receives scaled mass
\(nc_n(\theta)w\). The next proposition identifies the weak limit of these
two scaled edge measures before addressing their exact support sizes.
\begin{proposition}
\label{prop:atom-measure-limit}
Assume \(r_n\to r\in(0,1/4)\). If \(\nu_{n,+}\) and \(\nu_{n,-}\) are the
scaled edge measures of the NPMLE, then
\[
    (\nu_{n,+},\nu_{n,-})
    \Rightarrow
    (\nu_{b_r}^{*,+},\nu_{b_r}^{*,-}),
\]
where the two limits are independent copies of the optimizer in
Proposition~\ref{prop:atom-limit-unique}. In particular, their joint law
depends on the variance sequence only through \(r\).
\end{proposition}

\begin{proof}
Lemma~\ref{lem:support-fixed-core-collapse} isolates one central atom and
shows that the total physical edge mass tends to zero. The subthreshold
margin confines all remaining support from below. For every \(\delta>0\),
Lemma~\ref{lem:support-one-d-exclusion-margins} supplies a deterministic
\(c<1\) such that, except on an event of probability at most
\(\delta+o(1)\), both edge measures are supported on a fixed interval
\([a,c]\) containing \(b_r\). Proposition
\ref{prop:support-one-d-compact-likelihood} makes their total scaled masses
tight.

On a Skorokhod coupling, apply Lemmas
\ref{lem:atom-wall-functional-convergence} and
\ref{lem:atom-hard-wall}. Every subsequential limit maximizes the two
hard-wall functionals. Proposition~\ref{prop:atom-limit-unique} identifies
these maximizers uniquely. Letting first \(n\to\infty\) and then
\(\delta\downarrow0\) proves the joint convergence. The wall is \(b_r\) by
\eqref{eq:atom-penalty-exponent}.
\end{proof}

Weak convergence alone does not preserve an exact atom count: atoms with
vanishing mass could disappear, and several atoms could merge. The missing
input is nondegeneracy of every limiting KKT contact.
\begin{proposition}
\label{prop:atom-morse}
For every fixed \(b\in(1/2,1)\), the following statements hold almost surely
for the optimizer \(\nu_b^*\) and its KKT field \(\Gamma_{\nu_b^*}\).
\begin{enumerate}
\item The zero set of \(\Gamma_{\nu_b^*}\) on \([b,1)\) is exactly
      \(\supp(\nu_b^*)\).
\item Every interior contact \(\theta\) is quadratic:
      \(\Gamma_{\nu_b^*}''(\theta)<0\).
\item At the lower boundary, either \(\Gamma_{\nu_b^*}(b)<0\), or
      \(b\in\supp(\nu_b^*)\) and
      \(\Gamma_{\nu_b^*}'(b)<0\).
\end{enumerate}
\end{proposition}

The strict inequalities have the natural constrained interpretation. An
interior contact is a local maximum of a nonpositive function, while a
contact at the left endpoint has a nonpositive inward derivative.
Appendix~\ref{app:poisson-transversality} uses the continuously distributed
negative Poisson points to rule out equality and inactive contacts.

\subsection{Convergence of the atom count}

This subsection completes the proof of Theorem~\ref{thm:atom-count-limit} by
passing from weak convergence of the edge measures to convergence of their
active sets. Recall from \eqref{eq:atom-count-intro} that
\(\mathsf K_r=1+K_{b_r}^++K_{b_r}^-\), where \(b_r=1-\sqrt r\). The central
one is the atom isolated by Lemma~\ref{lem:support-fixed-core-collapse}.

The nondegenerate contacts of Proposition~\ref{prop:atom-morse} persist
one-for-one under the finite-sample approximation, including a possible
contact at the moving hard wall.
\begin{proposition}
\label{prop:atom-count-convergence}
If \(r_n\to r\in(0,1/4)\), then
\[
    |\supp(\widehat G_n)|\Rightarrow\mathsf K_r.
\]
Moreover, \(\Pbb\{\mathsf K_r=1\}=p(r)\).
\end{proposition}

\begin{proof}
Proposition~\ref{prop:atom-measure-limit} and positivity of every limiting
mass force at least one finite-sample atom into each small, disjoint
neighborhood of a limiting atom. Lemma
\ref{lem:fixed-exponent-support-tightness} gives the compactness needed to
extract the corresponding KKT fields. For the reverse bound, work on the
event from Proposition~\ref{prop:atom-measure-limit} on which both signed
edge measures are supported in a fixed compact interval and have bounded
scaled mass. Its probability is at least \(1-\delta-o(1)\).

Away from the finite limiting contact set, Proposition
\ref{prop:atom-morse} gives a strict negative KKT margin and excludes support
points. Near an interior active contact, the limiting second derivative is
strictly negative. Two-derivative convergence therefore makes the
finite-sample KKT field strictly concave there. A strictly concave,
nonpositive function has at most one zero, while weak convergence of the
positive limiting mass supplies at least one support point. Hence exactly one
finite-sample atom lies near each interior limiting atom.

It remains to treat the hard wall. If the boundary is inactive, its strict
KKT margin and the nonnegative penalty exclude a crossover contact. Suppose
instead that \(\nu_{b_r}^*\) has an atom at \(b_r\). On a fixed two-sided
neighborhood of \(b_r\), write the normalized finite-sample KKT field as
\[
    \Gamma_n(\theta)=R_n(\theta)-\cP_n(\theta).
\]
The penalty-free remainder satisfies \(R_n'<-a<0\), while direct
differentiation gives
\[
    \frac{\cP_n'}{\cP_n}=-(1-\theta)u_n^2+O(1),
    \qquad
    \frac{\cP_n''}{\cP_n}
       =(1-\theta)^2u_n^4+O(u_n^2).
\]
Wherever \(\Gamma_n'\ge0\), the positive term \(-\cP_n'\) is at least
\(a\); the first display then gives \(\cP_n\gtrsim u_n^{-2}\), and the
second gives \(\cP_n''\gtrsim u_n^2\). Since \(R_n''=O_{\Pbb}(1)\), we have
\(\Gamma_n''<0\) there. Thus \(\Gamma_n'\) cannot cross from negative to
positive. Two distinct zeros of the nonpositive function \(\Gamma_n\) would
force such a crossing, so there is at most one crossover contact. Weak
convergence of the boundary mass supplies at least one.

The argument applies independently at the two sample edges. Letting
\(\delta\downarrow0\) proves count convergence. Finally,
\(\nu_b^*=0\) exactly when \(\sup_{b\le\theta<1}S(\theta)\le0\). The two
edges are independent, so \eqref{eq:p-intro} gives
\(\Pbb\{\mathsf K_r=1\}=p(r)\).
\end{proof}

Proposition~\ref{prop:atom-morse} and the implicit-function theorem also
control changes in the hard-wall location. Under a common Poisson
realization, the atom count is locally constant at each fixed \(b\), almost
surely; this yields total-variation rather than only weak continuity.
\begin{proposition}
\label{prop:atom-count-continuity}
The map \(b\mapsto\mathcal L(K_b)\) is continuous in total variation on
\((1/2,1)\). Consequently,
\(r\mapsto\mathcal L(\mathsf K_r)\) is continuous in total variation on
\([0,1/4)\).
\end{proposition}

\begin{proof}
Couple all optimizers with the same Poisson process and fix deterministic
\(b\). If \(b\notin\supp(\nu_b^*)\), Proposition
\ref{prop:atom-morse} gives a strict boundary inequality. The support lies a
positive distance to the right of \(b\), and the KKT field is strictly
negative on a neighborhood of \(b\). For every \(b'\) sufficiently close to
\(b\), the same measure is feasible and satisfies the KKT conditions on
\([b',1)\). Uniqueness gives \(\nu_{b'}^*=\nu_b^*\).

Suppose instead that
\[
    \nu_b^*=s_1\delta_b+\sum_{j=2}^ks_j\delta_{\theta_j},
    \qquad b<\theta_2<\cdots<\theta_k.
\]
Parameterize this boundary stratum by its \(k\) masses and \(k-1\) interior
locations. The objective is analytic in these variables and in \(b\). Its
Hessian in the atomic variables is negative definite. Indeed, for a tangent
vector \(v\), put
\[
    g_v(x)=\sum_{j=1}^ke^{\theta_jx}
       \{\dot s_j+s_jx\dot\theta_j\},
    \qquad \dot\theta_1=0.
\]
At the critical point,
\[
 v^{\mathsf T}Hv
 =-\int_{\mathbb R}
       \frac{g_v(x)^2}{\{1+F_{\nu_b^*}(x)\}^2}\,N(dx)
   +\sum_{j=2}^ks_j
       \Gamma_{\nu_b^*}''(\theta_j)\dot\theta_j^2.
\]
The second term is nonpositive. If the first vanishes, \(g_v\) vanishes at
every negative Poisson point. Lemma~\ref{lem:atom-exp-polynomial} and linear
independence of the confluent exponentials give \(v=0\). Thus the Hessian is
nonsingular.

The implicit-function theorem supplies a nearby critical \(k\)-atomic
measure on the corresponding boundary stratum, with positive masses.
Proposition~\ref{prop:atom-morse} gives a strict negative KKT margin away
from the contacts, strict quadratic maxima at the interior contacts, and a
strictly negative inward derivative at the boundary contact. These
inequalities persist as \(b\) varies, so the critical branch satisfies the
global KKT conditions and is the unique optimizer. Hence \(K_{b'}=K_b\) for
all sufficiently close \(b'\).

We have shown that if \(b_m\to b\), then
\(\mathbf1_{\{K_{b_m}\ne K_b\}}\to0\) almost surely. Dominated convergence
therefore gives
\[
 d_{\mathrm{TV}}\{\mathcal L(K_{b_m}),\mathcal L(K_b)\}
 \le\Pbb\{K_{b_m}\ne K_b\}\longrightarrow0.
\]
The same coupling handles the sum of two independent copies, and
\(r\mapsto b_r\) is continuous on \((0,1/4)\). At zero,
\(\Pbb\{\mathsf K_r\ne1\}=1-p(r)\to0\) by Proposition
\ref{prop:phase-curve-properties}, which proves total-variation continuity on
\([0,1/4)\).
\end{proof}

\begin{proof}[Proof of Theorem~\ref{thm:atom-count-limit}]
For \(r\in(0,1/4)\), the theorem is Proposition
\ref{prop:atom-count-convergence} together with Proposition
\ref{prop:atom-count-continuity}. If \(r_n\to0\), Theorem
\ref{thm:exponent-phase-diagram} gives
\(\Pbb\{|\supp(\widehat G_n)|=1\}\to1\), so the count converges to
\(\mathsf K_0=1\).
\end{proof}

\subsection{Divergence at the critical endpoint}

This subsection proves
Proposition~\ref{prop:atom-count-endpoint-divergence}. In one-sided
coordinates, write
\(b=1/2+\beta\). The logarithmic distance from the wall to the endpoint
\(1/2\) then grows like \(\log(1/\beta)\). A positive density of separated
score locations can each contribute order-one likelihood gain, whereas a
single atom can exploit only the largest score on this interval.

For \(h>0\), define the normalized endpoint score and its negative-tail
quadratic variation by
\[
    X(h)=\sqrt{2h}\,S(1/2+h),
    \qquad
    Q_h=\int_{-\infty}^0e^{(1+2h)x}\,N(dx).
\]
The KMT construction in the proof of Proposition
\ref{prop:upper-endpoint-persistence} supplies all the uniform information
needed below. We collect those consequences together with the Poisson strong
law.
\begin{lemma}
\label{lem:atom-endpoint-control}
On an extension of the probability space, for every \(0<h_0<1/4\) there is
a centered stationary Gaussian process \(G_{h_0}\) with covariance
\[
    \E G_{h_0}(s)G_{h_0}(t)
    =\operatorname{sech}\left(\frac{t-s}{2}\right)
\]
such that
\[
    \sup_{t\ge0}|X(h_0e^{-t})-G_{h_0}(t)|
    \le R\sqrt{h_0},
\]
where \(R<\infty\) almost surely and can be chosen simultaneously for all
\(h_0<1/4\). Moreover,
\[
    \sup_{0\le t\le L}G_{h_0}(t)
       =O_{\Pbb}\{\sqrt{\log(L+2)}\}
\]
uniformly in \(h_0\), and each fixed-spacing sequence
\((G_{h_0}(j\Delta))_{j\ge0}\) is stationary and ergodic. Finally,
\[
    2hQ_h\longrightarrow1\quad\text{almost surely as }h\downarrow0,
    \qquad
    \sup_{0<h\le h_0}\sqrt h\,X(h)^+\longrightarrow0
       \quad\text{in probability as }h_0\downarrow0.
\]
\end{lemma}

\begin{proof}
Under \(t=e^{-x}\), the negative Poisson process becomes a unit-rate Poisson
process \(\Pi\) on \([1,\infty)\), and
\[
    \int_{-\infty}^0e^{(1/2+h)x}\{N-\mu\}(dx)
    =\int_1^\infty t^{-1/2-h}\{\Pi(dt)-dt\}.
\]
The infinite-horizon KMT coupling in
\eqref{eq:poisson-kmt} replaces the centered counting process by Brownian
motion with an error bounded by an almost surely finite random multiple of
\(1+\log t\). Integration by parts makes the normalized error
\(O(R\sqrt{h_0})\), uniformly for \(0<h\le h_0\). The positive half-line
contains only finitely many Poisson points almost surely, and its contribution
to the normalized process is another almost surely finite random multiple of
\(\sqrt{h_0}\). The normalized Brownian integrals have covariance
\[
    \frac{2\sqrt{hk}}{h+k};
\]
putting \(h=h_0e^{-t}\) gives the stated sech covariance.

A unit-interval Gaussian maximal inequality and a union bound give the
maximum estimate. Since the covariance tends to zero at infinity, every
fixed-spacing sampled sequence is ergodic. The same maximum estimate and
Borel--Cantelli show that
\(\sup_{t\ge0}e^{-t/2}G_{h_0}(t)^+<\infty\) almost surely, with a law
independent of \(h_0\); the coupling then gives the second limiting display.

Finally, if \(A(t)=\Pi([1,t])\), integration by parts gives
\[
    Q_h=(1+2h)\int_1^\infty A(t)t^{-2-2h}\,dt.
\]
The Poisson strong law \(A(t)/t\to1\), followed by the usual Abelian
splitting at a fixed large cutoff, yields \(2hQ_h\to1\).
\end{proof}

Let
\[
    V_b=\cJ_{N,b}(\nu_b^*)
\]
be the unrestricted one-sided optimum. A logarithmically spaced collection
of small atoms obtains likelihood proportional to the length of the endpoint
window.
\begin{proposition}
\label{prop:atom-many-gain}
There is a deterministic \(c>0\) such that
\[
    \Pbb\left\{V_{1/2+\beta}\ge c\log(1/\beta)\right\}
    \longrightarrow1
    \qquad(\beta\downarrow0).
\]
\end{proposition}

\begin{proof}
Put \(L_0=\log(1/\beta)\), \(H_\beta=L_0^{-4}\), and
\(L_\beta=\log(H_\beta/\beta)\sim L_0\). Fix \(\Delta>0\), and take
\[
    h_j=H_\beta e^{-j\Delta},
    \qquad 0\le j<J_\beta,
    \qquad J_\beta=\lfloor L_\beta/\Delta\rfloor.
\]
By Lemma~\ref{lem:atom-endpoint-control}, this score vector is uniformly
\(o_{\Pbb}(1)\) from a sampled stationary sech process. Choose \(q>0\).
Ergodicity gives constants \(p,q>0\) such that, with probability tending to
one, at least \(pJ_\beta\) indices satisfy \(X(h_j)\ge q\).

For a small constant \(a>0\), place mass \(a\sqrt{2h_j}\) at
\(1/2+h_j\) for each favorable index. The linear gain is at least
\(apqJ_\beta\). Since \(h(y)\ge-y^2/2\), its nonlinear cost is at most one
half of \(\int F^2\,dN\). Dominate this integral by the competitor using all
grid points. On the negative half-line its expectation is
\[
 a^2\sum_{j,l<J_\beta}\frac{2\sqrt{h_jh_l}}{h_j+h_l}
 =a^2\sum_{j,l<J_\beta}
    \operatorname{sech}\left(\frac{(j-l)\Delta}{2}\right)
 \le C_\Delta a^2J_\beta.
\]
Because the full-grid transform at zero is \(O(a\sqrt{H_\beta})\), the
variance of this Poisson integral is at most
\(C_\Delta a^4H_\beta J_\beta\). Its positive-half-line contribution is
\(o_{\Pbb}(1)\), since that half-line contains finitely many points.
Consequently,
\[
    \cJ_{N,1/2+\beta}(\nu_\beta)
    \ge\{apq-C_\Delta a^2+o_{\Pbb}(1)\}J_\beta.
\]
Choose \(a\) so that the deterministic coefficient is positive.
\end{proof}

The sparse upper bound rests on an exact subadditivity identity. It avoids
having to cover the class of all sparse measures.
\begin{lemma}
\label{lem:atom-gain-subadditivity}
If \(\nu=\sum_{j=1}^Ds_j\delta_{\theta_j}\), then
\[
    \cJ_{N,b}(\nu)
    \le\sum_{j=1}^D\cJ_{N,b}(s_j\delta_{\theta_j}).
\]
\end{lemma}

\begin{proof}
For \(y,z\ge0\),
\[
 h(y+z)-h(y)-h(z)
 =\log\frac{1+y+z}{(1+y)(1+z)}\le0.
\]
Induction proves subadditivity of the nonlinear term, and the score term is
linear. \qedhere
\end{proof}

Optimizing a single atom costs only the square of the largest normalized
score. The Gaussian maximum over logarithmic time therefore contributes the
smaller factor \(\log\log(1/\beta)\).
\begin{proposition}
\label{prop:atom-one-gain}
Let
\[
    M_\beta=
    \sup_{\beta\le h<1/2}\sup_{s\ge0}
       \{\cJ_{N,1/2+\beta}(s\delta_{1/2+h})\}_+.
\]
Then
\[
    M_\beta=O_{\Pbb}\{\log\log(1/\beta)\}.
\]
\end{proposition}

\begin{proof}
For \(x<0\), the inequality
\(h(y)\le-y^2/\{2(1+y)\}\) gives
\begin{equation}
\label{eq:atom-one-gain}
 \cJ_{N,1/2+\beta}(s\delta_{1/2+h})
 \le sS(1/2+h)-\frac{s^2Q_h}{2(1+s)}.
\end{equation}
Fix \(\delta>0\). Lemma~\ref{lem:atom-endpoint-control} allows us to choose
deterministic \(h_0>0\) such that, with probability at least \(1-\delta\),
\[
    Q_h\ge\frac1{4h},
    \qquad \sqrt h\,X(h)^+\le\frac1{16}
    \qquad(0<h\le h_0).
\]
On this event \(S(1/2+h)\le Q_h/4\). The right side of
\eqref{eq:atom-one-gain} is nonpositive for \(s\ge1\); for \(0\le s\le1\),
quadratic optimization gives
\[
    \cJ_{N,1/2+\beta}(s\delta_{1/2+h})
    \le\frac{S(1/2+h)_+^2}{Q_h}\le2X(h)_+^2.
\]
Lemma~\ref{lem:atom-endpoint-control} bounds the Gaussian maximum over this
\(O\{\log(1/\beta)\}\)-length interval by
\(O_{\Pbb}(\sqrt{\log\log(1/\beta)})\).

On \(h_0\le h<1/2\), the best one-atom gain is almost surely finite and
independent of \(\beta\). Indeed, \(S(\theta)\to-\infty\) as
\(\theta\uparrow1\), so only a random compact interval matters. Uniform
coercivity in \(s\) follows from the negative-point block argument in the
proof of Proposition~\ref{prop:support-one-d-analytic-active-set}. Since
\(\delta\) was arbitrary, the result follows.
\end{proof}

We now compare the unrestricted logarithmic gain with the best gain available
to a sparse optimizer.
\begin{proof}[Proof of Proposition~\ref{prop:atom-count-endpoint-divergence}]
Put \(r=r_\delta=1/4-\delta\) and
\(\beta=b_r-1/2=1/2-\sqrt r\), so \(\beta\sim\delta\). Write
\(a_\delta=\log(1/\delta)/\log\log(1/\delta)\). Fix \(\zeta>0\).
Proposition~\ref{prop:atom-one-gain} allows us to choose \(C<\infty\) so
that
\[
    \limsup_{\beta\downarrow0}
    \Pbb\{M_\beta>C\log\log(1/\beta)\}\le\zeta.
\]
On the complementary event, if \(K_{b_r}\le\eta a_\delta\), then
Lemma~\ref{lem:atom-gain-subadditivity} gives
\[
    V_{b_r}\le K_{b_r}M_\beta
       \le \{C\eta+o(1)\}\log(1/\beta).
\]
Proposition~\ref{prop:atom-many-gain} supplies a deterministic \(c>0\) for
which \(V_{b_r}\ge c\log(1/\beta)\) with probability tending to one. Thus,
for every \(\eta<c/(2C)\),
\[
    \limsup_{\delta\downarrow0}
    \Pbb\{K_{b_r}\le\eta a_\delta\}\le\zeta.
\]
Because \(\zeta\) was arbitrary, the nested limit is zero for the one-sided
count. Finally, \(\mathsf K_r=1+K_{b_r}^++K_{b_r}^-\), and the event in the
proposition implies \(K_{b_r}^+\le\eta a_\delta\), proving the stated result.
\end{proof}
